
%%% Kawauchi 156:
\subsection{Zgodność}
Zgodność jest relacją równoważności na zbiorze węzłów, która prowadzi do nowej definicji węzłów plastrowych (patrz fakt~\ref{prp:concordant_iff_sum_slice}).
My przytaczamy jej definicję z pracy Gompfa \cite{gompf1986}:

\begin{definition}[zgodność]
\index{zgodność}%
\index{węzeł!zgodny|see {zgodność}}%
    Dwa węzły $K_0, K_1$ nazywamy (gładko) zgodnymi, jeżeli zbiór
    \begin{equation}
        K_0 \times \{0\} \cup K_1 \times \{1\}
    \end{equation}
    jest brzegiem pewnego pierścienia gładko zanurzonego w $S^3 \times I$.
\end{definition}

Kawauchi \cite[s. 156]{kawauchi1996} pisze \emph{,,Two knots (…) are knot cobordant (or concordant)''}, więc tak jak wielu innych autorów nie odróżnia węzłów kobordantnych od zgodnych.
Mamy zamiar zrobić dokładnie to samo: różnica między tymi terminami jest subtelna; węzły zgodne są też kobordantne, ale implikacja w drugą stronę nie zachodzi (wiemy o~tym z~tekstu Blanlœila ,,Cobordism and Concordance of Knots'') chyba, że pracuje się z węzłami sferycznmi, a tak jest w klasycznej teorii węzłów.
\index[persons]{Blanloeil, Vincent}%
% https://www.maths.ed.ac.uk/~v1ranick/papers/blanloeil
% Concordant knots are cobordant, but the converse is not true in general.
% "Cobordism and Concordance of Knots" by Vincent Blanlœil

Dlatego my będziemy zawsze pisać o węzłach zgodnych i nigdy o kobordantnych.

\begin{proposition}
\label{prp:concordant_iff_sum_slice}%
    Dwa węzły $K_1, K_2$ są zgodne wtedy i tylko wtedy, gdy suma $(\operatorname{mr} K_2) \shrap K_1$ jest plastrowa.
\index{suma spójna}%
\index{węzeł!plastrowy}%
\end{proposition}

\begin{proof}
    Ćwiczenie 12.1.3 w książce Kawauchiego \cite{kawauchi1996}.
\end{proof}

\begin{definition}
    Węzeł zgodny z~niewęzłem nazywamy plastrowym.
\index{węzeł!plastrowy}%
\end{definition}

,,Bycie zgodnym'' jest relacją równoważności, słabszą od ,,bycia izotopijnym'', ale chyba mocniejszą od ,,bycia homotopijnym''.
\index{izotopia}%
\index{homotopia}
% ale mocniejszą od homotopii?
% izotopia: https://encyclopediaofmath.org/wiki/Cobordism_of_knots
% homotopia: https://en.wikipedia.org/wiki/Link_concordance By its nature, link concordance is an equivalence relation. It is weaker than isotopy, and stronger than homotopy: isotopy implies concordance implies homotopy. A link is a slice link if it is concordant to the unlink.
Klasę abstrakcji węzła $K$ oznaczamy przez $[K]$.

\begin{definition}[grupa zgodności]
\index{grupa!zgodności}%
    Niech $C^1$ oznacza iloraz zbioru wszystkich węzłów przez relację zgodności.
    Zbiór $C^1$ wyposażony w~działanie
    \begin{equation}
        [K_1] + [K_2] = [K_1 \shrap K_2]
    \end{equation}
    staje się grupą abelową, nazywaną grupą zgodności.
    Jej elementem neutralnym jest klasa abstrakcji niewęzła.
    Elementem przeciwnym do $[K]$ jest $[\operatorname{mr} K]$.
\end{definition}

%%% Kawauchi 157:

Niech $\Theta$ oznacza rodzinę macierzy Seiferta węzłów (czyli kwadratowych macierzy $V$ o~całkowitych wyrazach takich, że $\det (V - V^T) = 1$).
\index{macierz!Seiferta}%
Mówimy, że macierz $V \in \Theta$ jest zerowo kobordantna, jeżeli jest postaci
\index{macierz!Seiferta!zerowo kobordantna}%
\begin{equation}
    V = P \begin{pmatrix} 0 & V_{21} \\ V_{12} & V_{22} \end{pmatrix} P^{-1}
\end{equation}
dla pewnej całkowitoliczbowej macierzy $P$ o~wyznaczniku $\pm 1$; takie macierze nazywamy unimodularnie sprzężonymi.
\index{macierz!unimodularnie sprzężona}%
Każda zerowo kobordantna macierz $V \in \Theta$ stanowi macierz Seiferta pewnego plastrowego węzła $K$.
Kawauchi nazywa te węzły algebraicznie plastrowymi i~mówi, że to dokładnie węzły, które ograniczają izotropowe powierzchnie w kuli $B^4$, więc każdy węzeł plastrowy jest algebraicznie plastrowy.

Suma $(-V) \oplus V$ jest zerowo kobordantna dla każdej macierzy $V \in \Theta$.
To (chyba to) inspiruje Kawauchiego do wprowadzenia kolejnej definicji: dwie macierze $V_1, V_2 \in \Theta$ nazywa kobordantnymi, jeżeli $(-V_1) \oplus V_2$ jest zerowo kobordantna.
Kobordyzm stanowi relację równoważności na $\Theta$ -- iloraz $\Theta$ przez tę relację oznacza się $G_-$, jest grupą abelową.

\begin{proposition}
    % Kawauchi 12.2.8
    Odwzorowanie $\psi \colon C^1 \to G_-$ posyłające klasę abstrakcji węzła w klasę abstrakcji jego macierzy Seiferta jest dobrze określonym epimorfizmem.
\end{proposition}

\begin{proof}
    Nie umiemy nic udowodnić, więc wymienimy tylko trzy odsyłacze: z faktu~\ref{prp:cobordant_to_algebraic_is_algebraic} wynika, że odwzorowanie $\psi$ jest dobrze określone, dowód faktu~\ref{prp:signature_additive} pokazuje, że $\psi$ jest homomorfizmem, zaś w \cite[s. 62]{kawauchi1996} można przeczytać, dlaczego jest ,,na''.
\end{proof}

Funkcję $\psi$ rozpatrywał Levine \cite{levine1969} w latach sześćdziesiątych.
\index[persons]{Levine, Jerome}%
Po mniej niż dekadzie Casson, Gordon \cite{casson1978} wskazali nietrywialne elementy jądra.
\index[persons]{Casson, Andrew}%
\index[persons]{Gordon, Cameron}%
% to wyżej wiem z kawauchi98, "Supplementary notes for Chapter 12"
Potem był wynik Jianga \cite{jiang1981}, że jądro nie jest skończenie generowalne, bo zawiera izomorficzną kopię $\Z^\infty$, a~jeszcze później Livingstona \cite{livingston1999}, że zawiera też kopię $(\Z/2\Z)^\infty$.
% to wyżej wiem z https://mathscinet.ams.org/mathscinet-getitem?mr=2179265, pierwsze strony tekstu (nie recenzji)
\index[persons]{Jiang, Boju}%
\index[persons]{Livingston, Charles}%

\begin{proposition}
    $G_- \cong \Z^\infty \oplus (\Z/4\Z)^\infty \oplus (\Z/2\Z)^\infty$.
\end{proposition}

Kawauchi \cite[s. 161]{kawauchi1996} bez uzasadnienia postanawia nie przytoczyć dowodu tego faktu, ale opowiada krótko, jaka jest idea przewodnia i odsyła wprost do pracy Levine'a.
Na dalszych stronach jego pracy przeglądowej pojawiają się jakieś formy kwadratowe oraz uogólnienia wszystkiego do zgodności splotów, ale wracamy nocnym pociągiem i zaraz uśniemy...

\subsubsection{Nierozwiązane pytania o grupę zgodności}
O~strukturze samej grupy $C^1$ wiemy niewiele \cite[problem 1.38]{baykur2026}.
Fox i~Milnor \cite{fox1966} pokazali, że zawiera ona elementy rzędu dwa.
Przykładem jest ósemka $4_1$: jest równa swojemu lustrzanemu odwrotnemu, więc $[4_1] + [4_1] = 0$, ale nie jest plastrowa, bo jej wyznacznik wynosi $5$ (wniosek~\ref{det_slice_square}).
Hedden, Kim, Livingston \cite{hedden2016} znaleźli nieskończoną podgrupę złożoną z~węzłów topologicznie plastrowych, których klasy w~$C^1$ mają rząd dwa.
\index[persons]{Hedden, Matthew}%
\index[persons]{Kim, Se-Goo}%
Elementów nieskończonego rzędu jest dużo: grupa gładkich klas węzłów topologicznie plastrowych zawiera kopię $\Z^\infty$ (Endo \cite{endo1995}), a~nawet ma ją jako składnik prosty (Ozsváth, Stipsicz, Szabó \cite{ozsvath2017}; Dai, Hom, Stoffregen, Truong \cite{dai2021}).
\index[persons]{Endo, Hisaaki}%
\index[persons]{Stipsicz, András}%
\index[persons]{Truong, Linh}%

\begin{conjecture}
    \begin{enumerate}
        \item Czy w~$C^1$ istnieje element skończonego rzędu różnego od dwóch?
        Czy istnieje taki element w~podgrupie klas węzłów topologicznie plastrowych?
        \item Czy każdy element rzędu dwa jest zgodny z~pewnym silnie ujemnie zwierciadlanym węzłem (ang. \emph{strongly negative amphichiral})?
        \item Czy istnieje element nieskończenie podzielny, czyli klasa $[K]$ taka, że dla nieskończenie wielu liczb naturalnych $n$ istnieje węzeł $J$ spełniający $[K] = n [J]$?
    \end{enumerate}
\end{conjecture}

Element nieskończenie podzielny musi trafiać na zero pod każdym homomorfizmem $C^1 \to \Z$ (jak niezmienniki $\tau$ Ozsvátha-Szabó czy $s$ Rasmussena).

\paragraph{Hipoteza Rudolpha.}
Rudolph zapytał w~1976 roku, czy węzły algebraiczne, czyli węzły osobliwości krzywych płaskich (na przykład torusowe), są niezależne w~grupie zgodności, to znaczy generują jej wolną podgrupę \cite[problem 1.39]{baykur2026}.
Litherland \cite{litherland1979} udowodnił to dla węzłów torusowych, korzystając z~sygnatur Levine'a-Tristrama; jego dowód działa także dla zgodności topologicznej.
\index[persons]{Rudolph, Lee}%
\index[persons]{Litherland, Richard}%
\index{sygnatura!Levine'a-Tristrama}%
Ta metoda nie wystarcza w~ogólnym przypadku: Hedden, Kirk, Livingston \cite{hedden2012} pokazali, że istnieją sumy węzłów algebraicznych i~ich odbić, które są algebraicznie plastrowe (więc żadna sygnatura ich nie odróżni od niewęzła), a~mimo to nie są plastrowe.
\index[persons]{Kirk, Paul}%
Miyazaki \cite{miyazaki1994} dowiódł, że nietrywialne kombinacje liniowe iterowanych węzłów torusowych nie są taśmowe, zatem hipoteza plastrowo-taśmowa (patrz dalej) implikowałaby hipotezę Rudolpha.
\index[persons]{Miyazaki, Katura}%
Baker \cite{baker2016} oraz Abe, Tagami \cite{abe2016} zauważyli, że implikowałaby ona nawet jej mocniejszą wersję dla większej klasy pierwszych węzłów rozwłóknionych (silnie kwazidodatnich).
\index[persons]{Baker, Kenneth}%
\index[persons]{Abe, Tetsuya}%
\index[persons]{Tagami, Keiji}%
